\section{The \texttt{Virtual\_output} McStas Component}
Detector-like component that writes neutron state parameters into an ascci-format
'virtual source' neutron file.

\subsection*{Identification}
\begin{itemize}
  \item \textbf{Site:} 
  \item \textbf{Author:} \textless{}a href="mailto:farhi@ill.fr"\textgreater{}E. Farhi\textless{}/a\textgreater{}
  \item \textbf{Origin:} \textless{}a href="http://www.ill.fr"\textgreater{}ILL\textless{}/a\textgreater{}
  \item \textbf{Date:} Dec 17th, 2002
\end{itemize}

\subsection*{Description}
\begin{lstlisting}
Detector-like component writing neutron state parameters to a
virtual source neutron filename. The component geometry is the full
plane, and saves the neutron state as it exits from the previous
component.

It is particularly useful to generate a virtual source at a point that few
neutron reach. A long simulation will then only be performed once, to create
the upstream 'source' file. Further simulations are much faster if they start
from this low flux position with the 'source' filename.

The output file format is:
text column formatted with lines containing 11 values in the order:
p x y z vx vy vz t sx sy sz stored into about 83 bytes/n.

Beware the size of generated files ! When saving all events (bufsize=0) the
required memory has been optimized and remains very small. On the other hand
using large bufsize values (not recommended) requires huge storage memory.
Moreover, using the 'bufsize' parameter will often lead to wrong intensities.
Both methods will generate huge files.

A Vitess file may be obtained from the 'Vitess_output' component or from a
Vitess simulation (104 bytes per neutron) and read with Vitess_input.

Example: Virtual_output(filename="MySource.dat")
will generate a 9 Mo text file for 1e5 events stored.

%BUGS
Using bufsize non-zero may generate a virtual source with wrong intensity. This
component works with MPI (parallel execution mode).
\end{lstlisting}

\subsection*{Input parameters}
Parameters in \textbf{boldface} are required; the others are optional.

\begin{longtable}{p{0.22\textwidth}p{0.12\textwidth}p{0.46\textwidth}p{0.14\textwidth}}
\toprule
\textbf{Name} & \textbf{Unit} & \textbf{Description} & \textbf{Default} \\
\midrule
\endhead
filename & str & Name of neutron file to write. Default is standard output [string]. If not given, a unique name will be used. & 0 \\
bufsize & 1 & Size of neutron output buffer default is 0, i.e. save all - recommended. & 0 \\
\bottomrule
\end{longtable}

\subsection*{Links}
\begin{itemize}
  \item \href{run:/home/willend/willend-McCode/mcstas-comps/obsolete/Virtual_output.comp}{Source code} for \texttt{Virtual\_output.comp}.
\end{itemize}
\IfFileExists{Virtual_output_static.tex}{\input{Virtual_output_static.tex}}{}